\justifying
\NSFCBodyText

根据申请人已确认资料，以下在研项目与本课题存在方法能力或成果归属核对关系。本课题拟研究矿山影像变化定位和可靠性，各项目任务书中的具体重叠范围仍须逐项核对。

\textbf{（1）融合多源数据大模型驱动的光伏发电功率预测方法研究。}自治区高校基本科研业务费科研项目培育类，批准号XJEDU2026P005，经费5万元，执行期2026年1月至12月，申请人主持，负责总体方案、模型研究和成果组织。其对象为光伏功率预测，与本课题双时相影像的空间变化定位不同；通用表征、评测脚本及成果使用范围需结合任务书核对。

\textbf{（2）多模态融合技术研究及其在新疆产业集群中的应用。}自治区天池英才引进计划青年博士人才项目，经费50万元，执行期2025年2月至2028年2月，申请人主持，负责总体研究方案、关键模型设计和成果组织。依据已确认任务书，该项目无独立项目批准号。其覆盖范围较广，与本课题之间的数据、方法和考核成果边界需重点核对，不能仅依据题目认定完全不重叠。

\textbf{（3）面向政企数字化应用的智能知识库与问答系统研发。}新疆数字云创有限公司委托的横向课题，编号202604140013，经费12万元，执行期2026年6月至10月，申请人主持，按合同构建知识库，成果及知识产权由合同双方共同享有。本课题不以知识库问答为研究任务，若拟使用相关代码，须先核对合同授权与成果归属。

\textbf{（4）数据驱动下基于互补型生成对抗网络与神经统计模型的公路沙害智能预警研究。}国家自然科学基金地区项目，批准号62466056，经费32万元，执行期2025年1月至2028年12月，申请人作为参与者负责数据收集整理，基金委登记单位为新疆大学。其环境数据整理经验与本课题的数据核验有关，不能写为申请人主持的矿山研究基础。

\textbf{（5）可学习乘积量化与状态空间模型协同的长上下文高效注意力机制研究。}中国人工智能学会（CAAI）—蚂蚁科研基金，编号202604140006，经费15万元，执行期2026年2月至2027年2月，申请人参与，负责高效注意力机制研究。该项目与本课题具有通用模型方法的联系，本课题不将长上下文注意力作为核心任务。

上述项目与本课题的任务条款、代码授权、考核成果及列示范围，连同实际主要参与者的相关在研项目，将依据\textbf{【待补 F-MCT-03：任务书或合同对照及相关项目完整清单】}核定。拟通过逐项区分研究对象、数据、方法检验和成果归属明确边界，不在尚未完成对照时作无重复承诺。
